% Options for packages loaded elsewhere
\PassOptionsToPackage{unicode}{hyperref}
\PassOptionsToPackage{hyphens}{url}
%
\documentclass[
]{article}
\usepackage{amsmath,amssymb}
\usepackage{iftex}
\ifPDFTeX
  \usepackage[T1]{fontenc}
  \usepackage[utf8]{inputenc}
  \usepackage{textcomp} % provide euro and other symbols
\else % if luatex or xetex
  \usepackage{unicode-math} % this also loads fontspec
  \defaultfontfeatures{Scale=MatchLowercase}
  \defaultfontfeatures[\rmfamily]{Ligatures=TeX,Scale=1}
\fi
\usepackage{lmodern}
\ifPDFTeX\else
  % xetex/luatex font selection
\fi
% Use upquote if available, for straight quotes in verbatim environments
\IfFileExists{upquote.sty}{\usepackage{upquote}}{}
\IfFileExists{microtype.sty}{% use microtype if available
  \usepackage[]{microtype}
  \UseMicrotypeSet[protrusion]{basicmath} % disable protrusion for tt fonts
}{}
\makeatletter
\@ifundefined{KOMAClassName}{% if non-KOMA class
  \IfFileExists{parskip.sty}{%
    \usepackage{parskip}
  }{% else
    \setlength{\parindent}{0pt}
    \setlength{\parskip}{6pt plus 2pt minus 1pt}}
}{% if KOMA class
  \KOMAoptions{parskip=half}}
\makeatother
\usepackage{xcolor}
\usepackage[margin=1in]{geometry}
\usepackage{color}
\usepackage{fancyvrb}
\newcommand{\VerbBar}{|}
\newcommand{\VERB}{\Verb[commandchars=\\\{\}]}
\DefineVerbatimEnvironment{Highlighting}{Verbatim}{commandchars=\\\{\}}
% Add ',fontsize=\small' for more characters per line
\usepackage{framed}
\definecolor{shadecolor}{RGB}{248,248,248}
\newenvironment{Shaded}{\begin{snugshade}}{\end{snugshade}}
\newcommand{\AlertTok}[1]{\textcolor[rgb]{0.94,0.16,0.16}{#1}}
\newcommand{\AnnotationTok}[1]{\textcolor[rgb]{0.56,0.35,0.01}{\textbf{\textit{#1}}}}
\newcommand{\AttributeTok}[1]{\textcolor[rgb]{0.13,0.29,0.53}{#1}}
\newcommand{\BaseNTok}[1]{\textcolor[rgb]{0.00,0.00,0.81}{#1}}
\newcommand{\BuiltInTok}[1]{#1}
\newcommand{\CharTok}[1]{\textcolor[rgb]{0.31,0.60,0.02}{#1}}
\newcommand{\CommentTok}[1]{\textcolor[rgb]{0.56,0.35,0.01}{\textit{#1}}}
\newcommand{\CommentVarTok}[1]{\textcolor[rgb]{0.56,0.35,0.01}{\textbf{\textit{#1}}}}
\newcommand{\ConstantTok}[1]{\textcolor[rgb]{0.56,0.35,0.01}{#1}}
\newcommand{\ControlFlowTok}[1]{\textcolor[rgb]{0.13,0.29,0.53}{\textbf{#1}}}
\newcommand{\DataTypeTok}[1]{\textcolor[rgb]{0.13,0.29,0.53}{#1}}
\newcommand{\DecValTok}[1]{\textcolor[rgb]{0.00,0.00,0.81}{#1}}
\newcommand{\DocumentationTok}[1]{\textcolor[rgb]{0.56,0.35,0.01}{\textbf{\textit{#1}}}}
\newcommand{\ErrorTok}[1]{\textcolor[rgb]{0.64,0.00,0.00}{\textbf{#1}}}
\newcommand{\ExtensionTok}[1]{#1}
\newcommand{\FloatTok}[1]{\textcolor[rgb]{0.00,0.00,0.81}{#1}}
\newcommand{\FunctionTok}[1]{\textcolor[rgb]{0.13,0.29,0.53}{\textbf{#1}}}
\newcommand{\ImportTok}[1]{#1}
\newcommand{\InformationTok}[1]{\textcolor[rgb]{0.56,0.35,0.01}{\textbf{\textit{#1}}}}
\newcommand{\KeywordTok}[1]{\textcolor[rgb]{0.13,0.29,0.53}{\textbf{#1}}}
\newcommand{\NormalTok}[1]{#1}
\newcommand{\OperatorTok}[1]{\textcolor[rgb]{0.81,0.36,0.00}{\textbf{#1}}}
\newcommand{\OtherTok}[1]{\textcolor[rgb]{0.56,0.35,0.01}{#1}}
\newcommand{\PreprocessorTok}[1]{\textcolor[rgb]{0.56,0.35,0.01}{\textit{#1}}}
\newcommand{\RegionMarkerTok}[1]{#1}
\newcommand{\SpecialCharTok}[1]{\textcolor[rgb]{0.81,0.36,0.00}{\textbf{#1}}}
\newcommand{\SpecialStringTok}[1]{\textcolor[rgb]{0.31,0.60,0.02}{#1}}
\newcommand{\StringTok}[1]{\textcolor[rgb]{0.31,0.60,0.02}{#1}}
\newcommand{\VariableTok}[1]{\textcolor[rgb]{0.00,0.00,0.00}{#1}}
\newcommand{\VerbatimStringTok}[1]{\textcolor[rgb]{0.31,0.60,0.02}{#1}}
\newcommand{\WarningTok}[1]{\textcolor[rgb]{0.56,0.35,0.01}{\textbf{\textit{#1}}}}
\usepackage{graphicx}
\makeatletter
\def\maxwidth{\ifdim\Gin@nat@width>\linewidth\linewidth\else\Gin@nat@width\fi}
\def\maxheight{\ifdim\Gin@nat@height>\textheight\textheight\else\Gin@nat@height\fi}
\makeatother
% Scale images if necessary, so that they will not overflow the page
% margins by default, and it is still possible to overwrite the defaults
% using explicit options in \includegraphics[width, height, ...]{}
\setkeys{Gin}{width=\maxwidth,height=\maxheight,keepaspectratio}
% Set default figure placement to htbp
\makeatletter
\def\fps@figure{htbp}
\makeatother
\setlength{\emergencystretch}{3em} % prevent overfull lines
\providecommand{\tightlist}{%
  \setlength{\itemsep}{0pt}\setlength{\parskip}{0pt}}
\setcounter{secnumdepth}{-\maxdimen} % remove section numbering
\usepackage{amsmath}
\usepackage{mathtools}
\usepackage{amsthm}
\usepackage{multirow}
\usepackage{booktabs}
\ifLuaTeX
  \usepackage{selnolig}  % disable illegal ligatures
\fi
\IfFileExists{bookmark.sty}{\usepackage{bookmark}}{\usepackage{hyperref}}
\IfFileExists{xurl.sty}{\usepackage{xurl}}{} % add URL line breaks if available
\urlstyle{same}
\hypersetup{
  pdftitle={Discrete Multivariate Analysis - 579 - HW \#8},
  pdfauthor={Alex Towell (atowell@siue.edu)},
  hidelinks,
  pdfcreator={LaTeX via pandoc}}

\title{Discrete Multivariate Analysis - 579 - HW \#8}
\author{Alex Towell
(\href{mailto:atowell@siue.edu}{\nolinkurl{atowell@siue.edu}})}
\date{}

\begin{document}
\maketitle

\newcommand{\var}{\operatorname{Var}}
\newcommand{\expect}{\operatorname{E}}
\newcommand{\corr}{\operatorname{Corr}}
\newcommand{\cov}{\operatorname{Cov}}
\newcommand{\ssr}{\operatorname{SSR}}
\newcommand{\se}{\operatorname{SE}}
\newcommand{\mat}[1]{\mathbf{#1}}
\newcommand{\RR}{\operatorname{RR}}
\newcommand{\eval}[2]{\left. #1 \right\vert_{#2}}

\hypertarget{problem-1}{%
\section{Problem 1}\label{problem-1}}

Consider an experiment on chlorophyll inheritance in maize. A genetic
theory predicts the ratio of green to yellow to be 3:1. In a sample of
\(n = 1103\) seedlings, \(n_1 = 854\) were green and \(n_2 = 249\) were
yellow.

\hypertarget{part-a}{%
\subsection{Part (a)}\label{part-a}}

\fbox{Compute the statistic $G^2$ for testing the proposed model.}

The likelihood ratio statistic is defined as \[
  G^2 = - 2 \log \Lambda = 2 \sum_{j=1}^{c} \log\left(\frac{n_j}{n \pi_{j 0}}\right).
\]

We use the following R code to compute the observed statistic \(G_0^2\).

\begin{Shaded}
\begin{Highlighting}[]
\CommentTok{\# below is code for testing a specified multinomial}
\CommentTok{\# data is entered as a vector of counts}
\NormalTok{obs }\OtherTok{=} \FunctionTok{c}\NormalTok{(}\DecValTok{854}\NormalTok{,}\DecValTok{249}\NormalTok{)}
\NormalTok{n }\OtherTok{=} \FunctionTok{sum}\NormalTok{(obs)}

\CommentTok{\# the null model is entered as a vector of hypothesized probabilities}
\CommentTok{\# expected counts under the null model are computed as n*prob}
\NormalTok{pi}\FloatTok{.0} \OtherTok{=} \FunctionTok{c}\NormalTok{(.}\DecValTok{75}\NormalTok{,.}\DecValTok{25}\NormalTok{)}
\NormalTok{m }\OtherTok{=}\NormalTok{ n}\SpecialCharTok{*}\NormalTok{pi}\FloatTok{.0}

\CommentTok{\# log{-}likelihood ratio test}
\NormalTok{G2 }\OtherTok{=} \DecValTok{2}\SpecialCharTok{*}\FunctionTok{sum}\NormalTok{(obs}\SpecialCharTok{*}\FunctionTok{log}\NormalTok{(obs}\SpecialCharTok{/}\NormalTok{m))}
\FunctionTok{print}\NormalTok{(G2)}
\end{Highlighting}
\end{Shaded}

\begin{verbatim}
## [1] 3.539017
\end{verbatim}

We see that the observed statistic is given by \(G_0^2 = 3.54\).

\hypertarget{part-b}{%
\subsection{Part (b)}\label{part-b}}

\fbox{Compute the $p$-value for the proposed model.}

The reference distribution is the chi-squared distribution with
\(c-1 = 1\) degrees of freedom.

The \(p\)-value with \(k\) degrees of freedom is defined as \[
  \text{$p$-value} = \operatorname{P}(\chi^2(k) > G^2),
\] or in this case \[
  \text{$p$-value} = \operatorname{P}(\chi^2(1) > 3.54) = 0.06.
\] The following R code computes the \(p\)-value.

\begin{Shaded}
\begin{Highlighting}[]
\CommentTok{\# the p{-}value, a calculation based on tail probabilities}
\NormalTok{c }\OtherTok{=} \FunctionTok{length}\NormalTok{(obs)}
\NormalTok{pvalue.G2 }\OtherTok{=} \FunctionTok{pchisq}\NormalTok{(G2,c}\DecValTok{{-}1}\NormalTok{,}\AttributeTok{lower.tail=}\ConstantTok{FALSE}\NormalTok{)}
\FunctionTok{print}\NormalTok{(pvalue.G2)}
\end{Highlighting}
\end{Shaded}

\begin{verbatim}
## [1] 0.05994099
\end{verbatim}

\hypertarget{part-c}{%
\subsection{Part (c)}\label{part-c}}

\fbox{Provide an interpretation of your result, stated in the context of the problem.}

According to the Fisher scale for interpreting \(p\)-values,
\(p\)-values of \(0.05\) and \(0.1\) denote respectively \emph{moderate}
and \emph{boderline} evidence against the null model.

Since we obtained a \(p\)-value of \(0.06\), the data provides
\emph{moderate} evidence against the genetic theory positing that the
ratio of green to yellow seedlings is 3:1.

\hypertarget{part-d}{%
\subsection{Part (d)}\label{part-d}}

\fbox{What are the shortcomings of using a $p$-value as a measure of evidence?}

The \(p\)-value overstates evidence against the null model by comparing
the null model to the alternative model that is best supported by the
data.

\hypertarget{part-e}{%
\subsection{Part (e)}\label{part-e}}

\fbox{
\begin{minipage}{0.8\textwidth}
Compute the Bayes Factor $B_{0 1}$ for medium, wide, and ultrawide prior
distributions.
According to the Bayes Factors, which model is supported by the data?
\end{minipage}
}

Below is code for computing Bayes Factors in a test for a specified
probability.

\begin{Shaded}
\begin{Highlighting}[]
\FunctionTok{library}\NormalTok{(BayesFactor)}
\CommentTok{\# g denotes number of green seedlings}
\CommentTok{\# y denotes number of yellow seedlings}
\CommentTok{\# n = g+y denotes total number of seedlings}
\NormalTok{g }\OtherTok{=} \DecValTok{854}
\NormalTok{y }\OtherTok{=} \DecValTok{249}
\NormalTok{n }\OtherTok{=}\NormalTok{ g}\SpecialCharTok{+}\NormalTok{y}
\NormalTok{p0 }\OtherTok{=} \DecValTok{3}\SpecialCharTok{/}\DecValTok{4}

\NormalTok{bf.medium }\OtherTok{=} \DecValTok{1}\SpecialCharTok{/}\FunctionTok{proportionBF}\NormalTok{(}\AttributeTok{y=}\NormalTok{g,}\AttributeTok{N=}\NormalTok{n,}\AttributeTok{p=}\NormalTok{p0,}\AttributeTok{rscale=}\StringTok{"medium"}\NormalTok{)}
\NormalTok{bf.wide }\OtherTok{=} \DecValTok{1}\SpecialCharTok{/}\FunctionTok{proportionBF}\NormalTok{(}\AttributeTok{y=}\NormalTok{g,}\AttributeTok{N=}\NormalTok{n,}\AttributeTok{p=}\NormalTok{p0,}\AttributeTok{rscale=}\StringTok{"wide"}\NormalTok{)}
\NormalTok{bf.ultra }\OtherTok{=} \DecValTok{1}\SpecialCharTok{/}\FunctionTok{proportionBF}\NormalTok{(}\AttributeTok{y=}\NormalTok{g,}\AttributeTok{N=}\NormalTok{n,}\AttributeTok{p=}\NormalTok{p0,}\AttributeTok{rscale=}\StringTok{"ultrawide"}\NormalTok{)}

\NormalTok{bf.medium}
\end{Highlighting}
\end{Shaded}

\begin{verbatim}
## Bayes factor analysis
## --------------
## [1] Null, p=0.75 : 1.931436 ±0%
## 
## Against denominator:
##   Alternative, p0 = 0.75, r = 0.5, p =/= p0 
## ---
## Bayes factor type: BFproportion, logistic
\end{verbatim}

\begin{Shaded}
\begin{Highlighting}[]
\NormalTok{bf.wide}
\end{Highlighting}
\end{Shaded}

\begin{verbatim}
## Bayes factor analysis
## --------------
## [1] Null, p=0.75 : 2.700316 ±0%
## 
## Against denominator:
##   Alternative, p0 = 0.75, r = 0.707106781186548, p =/= p0 
## ---
## Bayes factor type: BFproportion, logistic
\end{verbatim}

\begin{Shaded}
\begin{Highlighting}[]
\NormalTok{bf.ultra}
\end{Highlighting}
\end{Shaded}

\begin{verbatim}
## Bayes factor analysis
## --------------
## [1] Null, p=0.75 : 3.796726 ±0.01%
## 
## Against denominator:
##   Alternative, p0 = 0.75, r = 1, p =/= p0 
## ---
## Bayes factor type: BFproportion, logistic
\end{verbatim}

If \(\rm{BF}_{0 1} > 1\), then the data supports the null model over the
alternative model.

According to the Bayes Factors, the results for \emph{medium},
\emph{wide} and \emph{ultrawide} are given respectively by \(1.9\),
\(2.7\), and \(3.8\). Thus, the null model is supported by the data in
each case.

\end{document}
